\chapter{Mapping to orGUI models and validation}
\label{chap:mapping}

\section{Flattened atomic representation}

The mathematical kernel should consume a flattened list of generated atomic
records rather than special-case each high-level model.  After stacking, the
records have the following meanings.

\begin{longtable}{@{}p{0.23\textwidth}p{0.67\textwidth}@{}}
\toprule
Model & Records supplied to the $F_{\bm h}$ sum \\
\midrule
\endfirsthead
\toprule
Model & Records supplied to the $F_{\bm h}$ sum \\
\midrule
\endhead
\texttt{UnitCell}
& Ordered structural-layer views, atoms, source lateral area, internal
coherent-domain transforms and occupancies, and the crystal-wide component
weight. \\
\texttt{Film}
& The finite cyclic layer cells generated by \texttt{createLayers}.  Their
domain translations place them at the physical stacked heights. \\
\texttt{EpitaxyInterface}
& Positive upper-material probability records; positive lower-material
records at their possibly strain-coupled positions; and negative lower-
material records that remove the sharp, unstrained bulk already present in
the terminal lattice sum. \\
\texttt{PoissonSurface}
& Positive records for grown material and negative records for material
removed from the sharp Film termination.  Termination corrections remain
coherent with the Film amplitude. \\
\bottomrule
\end{longtable}

Negative occupancies are amplitudes, not negative intensities.  They must pass
unchanged through Eqs.~\eqref{eq:atomic-channel} and
\eqref{eq:natural-density-share}.  Likewise, alternative crystal domains are
applied as coherent amplitudes within each configured domain; intensities are
combined only at the model level where the existing API explicitly defines an
incoherent mixture.

Water and other continuous-density models do not have an atomic UnitCell sum
and are outside this chapter.  They may use the same slab integral in
Eq.~\eqref{eq:profile-reference-term}, but their actual-density integral must
be derived from their own continuous profile.

\section{Preparation and evaluation boundary}

The prepared object must freeze all quantities that define the reference
problem:
\begin{itemize}
  \item the physical optical boundaries and the final, possibly simplified,
        values $\delta_m,\beta_m$;
  \item incident and reciprocal exit $k_z$, $A^\pm$, $A^\pm/k_z$, and reference
        planes for every medium;
  \item incident and exit polarization, external geometry, and the rod
        $\bm h=\Qpar$;
  \item the reference-area convention and the per-record reference shares
        satisfying Eq.~\eqref{eq:reference-partition}.
\end{itemize}
Atomic positions, occupancies, Debye--Waller parameters, and species may be
reevaluated against a frozen reference if that is the intended fit operation.
Changing lattice transforms, stacking geometry, photon energy, optical
boundaries, profile simplification, or the reference area invalidates the
preparation.  Rebuilding only the atomic list while silently changing the
optical profile would combine two different reference problems.

For each scattering point, evaluation is naturally ordered as follows:
\begin{enumerate}
  \item identify the rod from the Laue condition, Eq.~\eqref{eq:laue}, and
        calculate $\bm Q_m^\nu$, $P_{m,\nu}$, and $\Cchan_{m,\nu}$ for all
        media and all four channels;
  \item assign every transformed atom to its physical optical medium;
  \item accumulate Eq.~\eqref{eq:atomic-channel} coherently over atoms,
        domains, media, and channels, giving $F_{\bm h,u}^{\rm at}$;
  \item if $\bm h=\bm 0$, add the constant-profile contribution of
        Eq.~\eqref{eq:profile-reference-term}, completing $F_{\bm 0}$;
        otherwise $F_{\bm h}$ is already complete;
  \item convert to a rod amplitude with Eq.~\eqref{eq:rod-amplitude}, and only
        then form an intensity or combine with the prepared unperturbed
        reflection amplitude.
\end{enumerate}

\section{Observables}

Equation~\eqref{eq:rod-amplitude} already gives the amplitude of the emergent
plane wave on rod $\bm h$,
\[
  r_{\bm h}=-\frac{2\pi\ii r_e}{A_{\rm ref}k_{z,0,f}}F_{\bm h},
\]
normalised to one reference lateral cell, for every rod and every scattering
geometry.  Nothing further needs to be assumed about the measurement geometry:
fixed $\alpha_i$, fixed $\alpha_f$, $\alpha_i=\alpha_f$, coplanar or not, the
scan mode only decides which $(\alpha_i,\alpha_f)$ pair --- and hence which
$k_{z,m}$ and which $A^\pm$ --- are evaluated at each measured point.

On a genuine crystal truncation rod, $\bm h\neq\bm 0$, there is no unperturbed
beam in that direction and the observable follows from $r_{\bm h}$ alone.  On
the specular rod the scattered field is combined coherently with the prepared
plane-wave reflection amplitude,
\begin{equation}
  r_{\rm total}=r_0+r_{\bm 0},
  \qquad
  R_{\rm coh}=|r_{\rm total}|^2.
  \label{eq:specular-reflectivity}
\end{equation}
The $r_0$ in Eq.~\eqref{eq:specular-reflectivity} must be the same $r_0$ that
appears in Eq.~\eqref{eq:u-minus}, that is, it must come from the same prepared
incident field as the $A_{m,i}^{\pm}$ used in the matrix element.  For
unpolarized radiation, calculate $s$ and $p$ independently and average their
intensities; do not add their amplitudes.

\paragraph{Field intensity and the experimental input.}
The fitting design uses $R=|r|^2$ as a dimensionless field-intensity ratio,
including off-specular rod amplitudes.  Experimental data must already be
corrected to this convention and absolutely normalized when the fitted scale
is fixed to one.  For plane waves in vacuum, the ratio of normal powers
through a surface parallel to the sample is instead
$(\kappa_f/\kappa_i)|r|^2$, where
$\kappa_{i,f}=\kzero\sin\alpha_{i,f}>0$.  These two observables coincide in
specular reflection.  The normal-power ratio is not a universal conversion
from detector counts: illuminated area, footprint, and detector acceptance
belong to experimental reduction.  The fitter therefore adds no such ratio
and does not redefine $r$.  This follows the distinction between field
intensity, flux conservation, and differential scattering cross sections in
Renaud et al., Secs.~5.2.6 and 5.3.3--5.3.6.  Within each polarization pair,
the four DWBA propagation branches remain a coherent amplitude sum;
incoherent polarization mixtures are combined only after squaring.

\section{Returned quantities}

The kernel returns one matrix element and everything else is derived from it
by the prefactor of Eq.~\eqref{eq:rod-amplitude}.  Implementations write that
prefactor with the positive vacuum exit wavevector
$\kappa_f=\kzero\sin\alpha_f=-k_{z,0,f}$, so the sign appears reversed relative
to Eq.~\eqref{eq:rod-amplitude}; the two are identical.

\begin{center}
\begin{tabular}{@{}lll@{}}
\toprule
Quantity & Definition & Units \\
\midrule
$F_{\bm h}$
  & Eq.~\eqref{eq:Fh-definition}
  & electrons per $A_{\rm ref}$ \\
$r_{\bm h}$
  & $2\pi\ii r_e F_{\bm h}/(A_{\rm ref}\kappa_f)$
  & dimensionless \\
$r_0$
  & Fresnel amplitude of the prepared reference, Eq.~\eqref{eq:u-minus}
  & dimensionless \\
$r$
  & $r_0+r_{\bm h}$, with $r_0=0$ for $\bm h\neq\bm 0$
  & dimensionless \\
$F_{\rm eff}$
  & $A_{\rm ref}\kappa_f\,r/(2\pi\ii r_e)$
  & electrons per $A_{\rm ref}$ \\
$R$
  & $|r|^2$
  & dimensionless \\
\bottomrule
\end{tabular}
\end{center}

Three points are worth stating explicitly, because each has caused confusion.

\paragraph{Where $r_e$ belongs.}  The classical electron radius enters once,
inside the prefactor.  The separate convention, in which a bare matrix element
is turned into a cross-section kernel by $r_e^2|F_{\bm h}|^2$, applies to
$F_{\bm h}$ and not to $r$.  The two must never be mixed.

\paragraph{What $F_{\rm eff}$ is.}  Inverting the prefactor on the
\emph{total} amplitude gives the structure factor that a kinematic analysis
would infer from this reflectivity.  It is a DWBA output; it is not the
structure factor of a separate kinematical model, which the tutorial material
writes $F_{\rm kin}$.  The two are compared against each other, never
substituted for one another.  It is not linear in the model density,
because $r_0$ is zeroth order and exact in the reference to all orders, so it
is not the DWBA structure factor of the sample; that is $F_{\bm h}$.  It
vanishes as $\alpha_f\to0$ and is meaningful only well above the critical
angle.  Off-specular, where $r_0=0$, it degenerates to $F_{\bm h}$.  Its value
is the large-$Q_z$ limit, where $r_0\to4\pi r_e\overline\rho_f/Q_z^2$ and
$F_{\rm eff}\to F_{\bm h}+A_{\rm ref}\overline\rho_f/(\ii Q_z)$, that is, the
full kinematic structure factor of the actual density --- exactly the quantity
a kinematically reduced data set contains.

\paragraph{Polarization-reduced fitting magnitudes.}
Reserve $P$ for the conventional dimensionless polarization intensity factor,
and denote the amplitude prefactor by
$c=2\pi\ii r_e/(A_{\rm ref}\kappa_f)$.  If experimental structure-factor
data were polarization corrected, intensity was multiplied by $1/P$ and
amplitude magnitude by $1/\sqrt{P}$.  Public fitting metadata stores $P$;
an absent factor means no correction, equivalently $P=1$.  Let $\overline R$
be the resolution-broadened field intensity, with any incident mixture and
unanalysed outgoing states summed incoherently.  The predicted magnitude is
\[
  F_{\rm pred}=\frac{\sqrt{\overline R/P}}{|c|},
\]
using central-geometry $P$ and $c$.  There is no extra $\kappa_i$ in this
conversion, although the incident geometry still enters the wavefields.
For a narrow integrated region where $P$ changes little, a central factor
also approximates pixelwise polarization correction.  This is not an exact
identity for general integration weights; it must be adequate for the
experimental uncertainty.  It does not replace the polarization-dependent
DWBA fields or imply a unique complex $F_{\rm eff}$ for a mixed measurement.

\paragraph{Why $R=|r|^2$ is the reflectivity to use.}  On the specular rod at
small angles the problem is genuinely one-dimensional, so the scalar reduction
of Eq.~\eqref{eq:helmholtz} involves no approximation and the $s$ and $p$
channels become degenerate; and $r_0$ carries the exact Fresnel and refraction
response of the prepared reference to all orders.  The only error is therefore
second order in the density contrast, and it shrinks with decreasing angle:
below the critical angle the exact reflectivity is unity and $|r_0+r_{\bm h}|^2$
reproduces it.  A reference whose density is wrong by $10\%$ gives a relative
error near $10^{-4}$ at one tenth of the critical angle, rising to a few
percent near $\alpha_c$, and falling quadratically as the contrast error is
reduced.  Squaring the amplitude rather than expanding it also keeps $R$
non-negative.  The strictly linearised variant,
$|r_0|^2+2\operatorname{Re}(r_0^*r_{\bm h})$, is useful only as a diagnostic:
a large difference between the two means $|r_{\bm h}|$ is not small against
$|r_0|$ and the first-order treatment is being pushed beyond its regime.

\section{Regression invariants}

The following checks are required before the atomistic extension is considered
scientifically usable.

\begin{enumerate}
  \item \textbf{Four-channel coherence.}  A finite layer with nonzero
        $A_i^\pm$ and $A_f^\pm$ equals the direct atom-by-atom evaluation of
        Eq.~\eqref{eq:Fh-definition}; changing to a sum of channel intensities
        must fail the test.
  \item \textbf{Born limit.}  Setting $A_i^+=A_f^+=1$ and the other branches to
        zero recovers the conventional UnitCell structure factor and the
        standard polarization contraction.
  \item \textbf{Born--Fresnel amplitude.}  A uniform half-space against a
        vacuum reference reproduces Eq.~\eqref{eq:born-fresnel}, including its
        sign.  This pins the prefactor of Eq.~\eqref{eq:rod-amplitude}
        independently of any external convention.
  \item \textbf{Terminal limit.}  A finite stack made progressively thicker
        converges to Eqs.~\eqref{eq:bulk-atomic-sum} and
        \eqref{eq:bulk-reference} away from exact Bragg poles.
  \item \textbf{Identical actual and reference density.}  Replacing every
        optical slab by its exactly matching atomic forward density gives
        $F_{\bm 0}=0$ within the atomic-truncation accuracy.
  \item \textbf{Off-specular invariance.}  Adding the reference-density
        bookkeeping does not change any $F_{\bm h}$ with $\bm h\neq\bm0$, as
        required by Eq.~\eqref{eq:coeff-nonspec}.
  \item \textbf{Partition invariance.}  Summing per-record results from
        Eq.~\eqref{eq:per-cell-specular} equals the profile-level
        Eq.~\eqref{eq:total-specular}, including after profile coalescing and
        simplification.
  \item \textbf{Area invariance.}  Two commensurate descriptions with different
        source lateral cells give the same amplitude per $A_{\rm ref}$ after
        Eq.~\eqref{eq:area-scale} is applied to both atomic and continuum
        terms.
  \item \textbf{Signed-model identity.}  At zero width, zero growth, or zero
        etching, the positive and negative interface/surface records cancel in
        both contributions to $F_{\bm h}$.
  \item \textbf{Boundary split.}  Moving an atom across an optical boundary
        changes its channel coefficient and $Q_z$ exactly once; a cell
        straddling a boundary agrees with the same atom list represented as two
        records.
  \item \textbf{Low-$Q_z$ stability.}  Equation~\eqref{eq:phi-finite} approaches
        the slab thickness without catastrophic cancellation.
  \item \textbf{Reciprocity.}  Exchanging source and detector geometry obeys
        the reciprocal-field symmetry of
        Eq.~\eqref{eq:u-minus-identification} without conjugating the exit
        field.
  \item \textbf{Wronskian normalisation.}  For any prepared stack, the
        Wronskian of Eq.~\eqref{eq:wronskian-value} evaluates to
        $2\ii k_{z,0}$ independently of the layer structure.  A solver whose
        $A^\pm$ violate this has an internal normalisation error that
        Eq.~\eqref{eq:rod-amplitude} would silently absorb.
  \item \textbf{No duplicate attenuation.}  Absorption carried by complex
        $Q_z$ is not also applied as an empirical depth factor.
\end{enumerate}

\section{References}

The electric-field, signed-$k_z$, reciprocal-exit-field, and four-channel
semantics follow G.~Renaud, R.~Lazzari, and F.~Leroy, ``Probing surface and
interface morphology with grazing incidence small angle X-ray scattering,''
\emph{Surface Science Reports} \textbf{64} (2009), 255--380, especially
Sections~5.2--5.4, Eq.~(91), and the four-event finite-layer examples in
Section~6.2.  The present note retains their electromagnetic derivation but
does not adopt the scalar wavefunction notation used for the optional quantum-
mechanical analogy.

Renaud's Eq.~(100) quotes the one-dimensional stratified Green function, with
its factor $1/k_z$, without derivation, citing P.~M.~Morse and H.~Feshbach,
\emph{Methods of Theoretical Physics} (McGraw--Hill, 1953) for the general
construction, together with X.~L.~Zhou and S.~H.~Chen, \emph{Physics Reports}
\textbf{257} (1995), 223--348; A.~V.~Andreev, A.~G.~Michette and A.~Renwick,
\emph{Journal of Modern Optics} \textbf{35} (1988), 1667; V.~F.~Sears,
\emph{Physical Review B} \textbf{48} (1993), 17477; I.~D.~Feranchuk,
\emph{Physical Review B} \textbf{52} (1995), 9214; and A.~A.~Maradudin and
D.~L.~Mills, \emph{Physical Review B} \textbf{11} (1975), 1392.
Section~\ref{sec:green} carries out that construction explicitly in orGUI's
own conventions, so the present document does not depend on any of them.

The terminal crystal-truncation result and the atomic-density truncation caveat
are cross-checked against V.~M.~Kaganer, ``Crystal truncation rods in
kinematical and dynamical x-ray diffraction theories,'' \emph{Physical Review
B} \textbf{75} (2007), 245425.  The perturbation as a difference from a planar
reference follows the same physical construction used by G.~H.~Vineyard,
\emph{Physical Review B} \textbf{26} (1982), 4146--4159, and by S.~K.~Sinha
\emph{et al.}, \emph{Physical Review B} \textbf{38} (1988), 2297--2311.

The companion file \texttt{dwba\_literature\_crosscheck.md} records the broader
comparison with the surface-DWBA and reflectivity literature.
